% Introduction -- hand-written.  Everything after this is generated from
% ../notes/*.md.

This book builds a neural-network library from nothing, one piece at a
time, and uses each piece on sound and music as soon as it exists.  The
library is called \texttt{weft} (\emph{Weights Estimated From Training}).
It is written in C++, it is header-only, and it depends on nothing beyond
the C++ standard library: the matrix product, backpropagation, the
optimisers, the Fourier transform, convolution and attention are all
implemented here and explained here.

The book and the library were written together.  Each chapter introduces
one part of \texttt{weft}: it states the mathematics, derives the
gradients, shows the code that implements them, and points to a test that
checks the result numerically.  Nothing in the library is there without a
chapter that explains it, and no chapter describes code that you cannot
open and run.

\section*{What the chapters cover}

The first chapters build the core: matrices, backpropagation and the dense
layer, activations, the cross-entropy loss, the network container,
optimisers and dropout.  Two small classifiers, on the Iris flowers and on
handwritten digits, show the core working end to end.

Audio enters with the fast Fourier transform and spectral features, which
lead to a classifier of orchestral instrument samples.  Autoencoders and
variational autoencoders follow, first on digits and then on spectrograms,
where they are used to generate sounds and to interpolate between timbres.

A chapter on performance makes training practical with a threaded matrix
product and an optional BLAS backend.  Four-dimensional tensors then open
the way to convolutional networks and to a convolutional autoencoder for
audio.  A chapter on differentiable synthesis then lets a network control
a classical harmonic-plus-noise synthesiser, which is where the book's
audio becomes clean.  The last chapters build attention, positional
encoding and layer normalisation, and assemble them into a transformer
block.

\section*{What you need}

You need to be comfortable reading C++ and to remember some linear algebra
and calculus: matrix products, partial derivatives and the chain rule.  No
earlier experience with machine learning is assumed.  The first chapter
explains why the library is in C++ and not in Python.

\section*{How to read it}

Read with the source open.  The headers are in \texttt{src/}, the tests in
\texttt{tests/} and the complete programs in \texttt{examples/}.  From the
\texttt{build/} directory, \texttt{make} compiles everything and
\texttt{make test} runs the tests.  The chapters are meant to be read in
order, because each one uses the code written in the chapters before it.
